\section{Introduction}

Quantitative trading seeks to turn noisy market observations into signals and portfolios that generate reliable out-of-sample returns. Formulaic alphas provide an executable and interpretable representation of such signals \citep{Kakushadze2016Formulaic}. A quantitative research workflow repeatedly formulates a market hypothesis, translates it into a computational strategy, evaluates the resulting signal, and uses the outcome to choose what to test next. Financial data are noisy and non-stationary, the candidate space is large, and only a fraction of hypotheses remain useful beyond the observations that motivated them. Large language models (LLMs) can support this workflow through hypothesis generation, code synthesis, and analysis of experimental feedback. Recent alpha-discovery methods illustrate the growing range of research procedures built around these capabilities \citep{Wang2023AlphaGPT,Tang2025AlphaAgent,Han2026QuantaAlpha}.

Evaluating an LLM research agent introduces three closely related challenges. First, a final strategy reflects both the underlying model and the harness that organizes its prompts, tools, memory, candidate revisions, and stopping decisions. Comparisons therefore require a clearly specified unit of evaluation. Second, historical market data support repeated adaptive experimentation, creating opportunities for selection bias and backtest overfitting \citep{Bailey2017Backtest}. The research feedback available to an agent must be separated from the evidence used to judge its frozen output. Third, operational choices affect the outcome: unequal model budgets, different data access, silent retries, and inconsistent return conventions can change the apparent ranking of otherwise comparable systems. A benchmark must make these choices part of the evaluation contract.

These challenges motivate a system-level question: how effectively can an LLM agent use a fixed research budget to produce predictive signals and cost-adjusted research returns under a fixed historical protocol? Answering this question requires observing the research process, including unsuccessful attempts, resource consumption, and final selection. It also requires precise financial semantics. A cross-sectional signal statistic, a fixed-notional research return, and a funded account outcome measure different properties of a strategy. Their interpretation depends on the information cutoff, holding period, missing data, and costs. In particular, annualizing the mean of a few observable cohorts does not establish performance over a fully observed calendar interval.

We introduce \textbf{\BenchmarkName}, a benchmark and execution platform for budgeted quantitative research on real historical market data. A participant consists of a registered language model and a research harness. The harness proposes executable stock-scoring strategies, receives authorized research feedback, and commits one final strategy per replicate for sealed assessment. One-shot generation, ordinary feedback iteration, and external harnesses can therefore be compared as explicit research procedures. Section~\ref{sec:benchmark} defines the suite that holds their evaluation conditions fixed.

The Arena platform makes the research process inspectable across service boundaries. Strategy programs access market observations through date-scoped authorization, while trusted services retain model credentials, evaluate outputs, and settle resource usage. Durable operation records allow recoverable interruptions to resume accepted work without silently generating a new model response. A shared financial core connects internal factor research and Arena assessment, with access and disclosure policies specific to each setting.

The daily benchmark developed here uses raw-price observations for a registered Shenzhen A-share population. It fixes a signal--entry--exit alignment, Top-$K$ selection with complete boundary ties, half-notional overlapping cohorts, and an explicit cost convention. Financial outcomes are accompanied by coverage, predictive statistics, and resource observations with completeness flags. Metric availability and ranking eligibility are reported separately. The current diversity gate can return zero when quality evidence is missing; its feature-level reasons are therefore essential to interpretation. This paper describes the protocol and platform implementation. Campaign results are left open.

The contributions of this work are:
\begin{itemize}
    \item \textbf{A benchmark for budgeted research episodes.} We define the model--harness pair as the participant and evaluate its research and selection process, retaining unsuccessful attempts and resource records.
    \item \textbf{An explicit financial evaluation protocol.} We formalize temporal alignment, portfolio selection, research returns, missing observations, and score qualification, with separate evidence for simulated account behavior.
    \item \textbf{An Arena execution platform.} We implement controlled market-data access, isolated strategy execution, durable budget accounting, final strategy commitment, and recoverable evaluation, together with campaign records for reproducible comparison.
\end{itemize}
